% =====================================================================
% TEMPLATE TUGAS AKHIR FMIPA UGM - DRAF STARTER PROJECT
% =====================================================================
% Seluruh bagian naskah ditulis dalam satu file ini.
%
% PENTING:
% 1. Template ini diadaptasi dari template asli menggunakan class yang diberikan pada panduan.
% 2. Pastikan file memilih cls yang sesuai.
% 3. Ubah data identitas pada bagian A di bawah ini sebelum kompilasi.
% 4. Kerangka Bab I-V di bawah merupakan kerangka awal. Cocokkan kembali
%    dengan Buku Panduan Penulisan Tugas Akhir FMIPA UGM dan arahan pembimbing.
% 5. Bab tutorial LaTeX diaktifkan untuk membantu pengguna mempelajari fitur.
%    Matikan dengan mengubah \showlatexguidetrue menjadi \showlatexguidefalse
%    sebelum menghasilkan naskah akhir untuk diserahkan.
%
% ---------------------------------------------------------------------
% PILIH SATU DOCUMENT CLASS
% ---------------------------------------------------------------------
% Skripsi:
\documentclass[skripsi]{ugmskripsi}
%
% Tesis:
%\documentclass[tesis]{ugmtesis}
%
% Disertasi:
% \documentclass[disertasi]{ugmdisertasi}
%
% Jangan mengaktifkan lebih dari satu perintah \documentclass.
% Opsi setiap class harus mengikuti deklarasi pada file .cls masing-masing.

% =====================================================================
% PAKET TAMBAHAN
% =====================================================================
% Class template telah memuat graphicx, natbib, dan nomencl.
% Karena itu, paket-paket tersebut tidak perlu dimuat ulang.

\usepackage{amsmath,amssymb} % Persamaan, aligned, matriks, dan simbol matematika.
\usepackage{booktabs}        % Garis tabel dengan kualitas tipografi yang baik.
\usepackage[hidelinks]{hyperref} % Tautan dan referensi silang pada PDF.

% ---------------------------------------------------------------------
% PENGATURAN BAB TUTORIAL LATEX
% ---------------------------------------------------------------------
% true  : bab panduan fitur LaTeX ditampilkan pada PDF.
% false : bab panduan fitur LaTeX tidak ditampilkan pada PDF.
\newif\ifshowlatexguide
\showlatexguidetrue

% =====================================================================
% A. DATA IDENTITAS TUGAS AKHIR
% =====================================================================
% Ganti seluruh teks contoh dengan data yang benar.
% Perhatikan kapitalisasi judul dan format identitas sesuai panduan Fakultas.

\titleind{TULISKAN JUDUL TUGAS AKHIR DALAM BAHASA INDONESIA}

\titleeng{WRITE THE THESIS TITLE IN ENGLISH}

\fullname{NAMA LENGKAP MAHASISWA}

\idnum{NOMOR INDUK MAHASISWA}

\examdate{Tanggal Bulan Tahun}

\degree{Gelar Akademik}

\yearsubmit{2026}

\program{Fisika}

\headprogram{Nama Ketua Program Studi}

\dept{Fisika}

\firstsupervisor{Nama Pembimbing Pertama}

\secondsupervisor{Nama Pembimbing Kedua}

\firstexaminer{Nama Penguji Pertama}

\secondexaminer{Nama Penguji Kedua}

\thirdexaminer{Nama Penguji Ketiga}

% =====================================================================
% MULAI DOKUMEN
% =====================================================================
\begin{document}

% =====================================================================
% B. BAGIAN AWAL TUGAS AKHIR
% =====================================================================
% Urutan perintah berikut mengikuti Template_Tesis.tex dari Fakultas.
% Jangan menghapus perintah khusus class tanpa memeriksa dampaknya.

% Sampul dan halaman judul
\cover
\titlepageind

% Halaman pengesahan dan pernyataan
\approvalpage
\declarepage

% ---------------------------------------------------------------------
% Halaman persembahan (opsional, sesuaikan dengan panduan Fakultas)
% ---------------------------------------------------------------------
\acknowledment
% Tuliskan persembahan secara singkat apabila bagian ini digunakan.
\begin{flushright}
\emph{Tuliskan kalimat persembahan di sini.}
\end{flushright}

% ---------------------------------------------------------------------
% Halaman motto (opsional, sesuaikan dengan panduan Fakultas)
% ---------------------------------------------------------------------
\motto
% Tuliskan motto atau kutipan yang ingin dicantumkan.
\emph{Tuliskan motto di sini.}

% ---------------------------------------------------------------------
% Prakata
% ---------------------------------------------------------------------
\preface
% Prakata berisi pengantar singkat mengenai penyusunan TUGAS AKHIR.
% Uraikan konteks penyusunan naskah dan ucapan terima kasih yang relevan.
% Hindari memasukkan pembahasan ilmiah yang seharusnya berada pada bab utama.

Tuliskan prakata TUGAS AKHIR di sini. Sesuaikan isi dan panjang prakata dengan
ketentuan pada panduan penulisan dan kebutuhan penulis.

\begin{flushright}
Yogyakarta, Tanggal Bulan Tahun\\
Penulis
\end{flushright}

% ---------------------------------------------------------------------
% Daftar isi dan daftar pendukung
% ---------------------------------------------------------------------
\tableofcontents
\listoftables
\listoffigures

% Jika daftar lambang dan singkatan digunakan, hapus tanda komentar
% pada perintah berikut. Entri lambang ditambahkan menggunakan \nomenclature.
% \lambang

% ---------------------------------------------------------------------
% Intisari Bahasa Indonesia
% ---------------------------------------------------------------------
\begin{abstractind}
% Intisari merangkum pokok penelitian secara ringkas.
% Jelaskan latar belakang atau konteks, tujuan, metode, hasil utama,
% dan simpulan. Ikuti batas panjang dan format kata kunci pada panduan.

Tuliskan intisari TUGAS AKHIR dalam bahasa Indonesia di sini.

\bigskip
Kata kunci: kata kunci pertama, kata kunci kedua, kata kunci ketiga.
\end{abstractind}

% ---------------------------------------------------------------------
% Abstract in English
% ---------------------------------------------------------------------
\begin{abstracteng}
% Abstract menyampaikan isi yang setara dengan intisari dalam bahasa Inggris.
% Pastikan istilah teknis, tujuan, metode, hasil, dan simpulan konsisten.

Write the thesis abstract in English here.

\bigskip
Keywords: first keyword, second keyword, third keyword.
\end{abstracteng}

% =====================================================================
% C. BAGIAN UTAMA TUGAS AKHIR
% =====================================================================
% Kerangka Bab I-VI di bawah adalah kerangka awal yang lazim digunakan.
% Susunan akhir, nama bab, dan subbab wajib disesuaikan dengan Buku Panduan
% Penulisan Tugas Akhir FMIPA UGM yang berlaku serta arahan pembimbing.

% ---------------------------------------------------------------------
% BAB I. PENDAHULUAN
% ---------------------------------------------------------------------
\chapter{PENDAHULUAN}

% Latar belakang menjelaskan konteks ilmiah, perkembangan topik,
% kesenjangan penelitian, dan alasan penelitian perlu dilakukan.
\section{Latar Belakang Masalah}
Tuliskan latar belakang penelitian di sini. Mulailah dari konteks umum,
kemudian arahkan pembahasan menuju masalah ilmiah yang lebih spesifik.
Jelaskan kesenjangan pengetahuan atau persoalan yang belum terselesaikan,
serta alasan penelitian ini penting untuk dilakukan.

% Rumusan masalah menyatakan pertanyaan utama yang akan dijawab penelitian.
\section{Perumusan Masalah}
Tuliskan permasalahan penelitian dalam bentuk uraian atau pertanyaan yang
jelas, spesifik, dan dapat dijawab melalui metode penelitian yang digunakan.

% Tujuan harus berhubungan langsung dengan rumusan masalah.
\section{Tujuan Penelitian}
Jelaskan hasil atau capaian ilmiah yang ingin diperoleh. Susun tujuan secara
spesifik dan selaras dengan setiap rumusan masalah.

% Manfaat menjelaskan kontribusi teoretis maupun praktis bila relevan.
\section{Manfaat Penelitian}
Jelaskan kontribusi penelitian terhadap pengembangan ilmu, metode, data,
atau penerapan yang relevan dengan bidang penelitian.

% Sistematika penulisan menjelaskan susunan bab dan subbab, serta alur logis
\section{Sistematika Penulisan}
Jelaskan susunan bab dan subbab dalam TUGAS AKHIR ini, serta alur logis dari pendahuluan hingga kesimpulan.
Uraikan bagaimana setiap bagian saling terkait dan mendukung satu sama lain.

% ---------------------------------------------------------------------
% BAB II. TINJAUAN PUSTAKA
% ---------------------------------------------------------------------
\chapter{TINJAUAN PUSTAKA}

% Tinjauan pustaka merangkum penelitian terdahulu yang relevan.
\section{Tinjauan Pustaka}
Bahas literatur dan hasil penelitian terdahulu yang berkaitan langsung dengan
topik TUGAS AKHIR. Susun pembahasan secara tematis atau kronologis, bukan sekadar
menyalin ringkasan setiap sumber. Tunjukkan hubungan literatur tersebut
dengan masalah penelitian yang dikaji.

% ---------------------------------------------------------------------
% BAB III. DASAR TEORI
% ---------------------------------------------------------------------
\chapter{DASAR TEORI}

% Dasar teori menyajikan konsep, prinsip, dan persamaan yang digunakan.
\section{Dasar Teori}
Uraikan konsep dan teori yang diperlukan untuk memahami serta menyelesaikan
permasalahan penelitian. Definisikan besaran, simbol, asumsi, dan notasi
secara konsisten. Sertakan penurunan persamaan apabila diperlukan untuk
menjelaskan metode atau model yang digunakan.

% Subbab berikut dapat ditambah sesuai kebutuhan bidang penelitian.
\subsection{Konsep atau Model Teoretis Pertama}
Jelaskan konsep, model, atau prinsip pertama yang digunakan dalam penelitian.

\subsection{Konsep atau Model Teoretis Kedua}
Jelaskan konsep, model, atau prinsip berikutnya yang diperlukan.

% Kerangka konsep atau hipotesis hanya digunakan jika relevan dengan desain.
\section{Kerangka Pemikiran dan Hipotesis}
Jelaskan hubungan antar-konsep atau alur logis yang menghubungkan teori,
masalah penelitian, dan pendekatan penyelesaian. Tuliskan hipotesis apabila
penelitian memang menggunakan hipotesis. Jika tidak relevan, bagian ini dapat
diganti atau dihilangkan sesuai panduan dan arahan pembimbing.

% ---------------------------------------------------------------------
% BAB IV. METODE PENELITIAN
% ---------------------------------------------------------------------
\chapter{METODE PENELITIAN}

% Jelaskan rancangan penelitian secara sistematis agar dapat ditelusuri.
\section{Rancangan dan Pendekatan Penelitian}
Jelaskan jenis penelitian dan pendekatan yang digunakan, misalnya teoretis,
eksperimental, komputasional, atau gabungan, sesuai karakter penelitian.

\section{Objek, Bahan, dan Peralatan Penelitian}
Jelaskan objek yang diteliti, bahan yang digunakan, serta perangkat atau
peralatan yang diperlukan. Untuk penelitian teoretis atau komputasional,
uraikan sistem fisika, data, perangkat lunak, dan sumber daya komputasi yang
relevan. Hilangkan unsur yang memang tidak digunakan.

\section{Variabel dan Parameter Penelitian}
Definisikan variabel, parameter, kondisi awal, kondisi batas, dan besaran
lain yang digunakan. Sertakan satuan serta rentang nilai apabila relevan.

\section{Prosedur atau Tahapan Penelitian}
Uraikan tahapan penelitian secara berurutan. Setiap tahapan perlu menjelaskan
apa yang dilakukan dan bagaimana tahapan tersebut menjawab tujuan penelitian.
Diagram alir dapat digunakan untuk membantu menjelaskan prosedur.

\section{Metode Pengolahan dan Analisis Data}
Jelaskan metode analisis, teknik perhitungan, prosedur validasi, serta kriteria
yang digunakan untuk menilai hasil. Tuliskan persamaan utama dan asumsi yang
mendasari analisis apabila diperlukan.

% ---------------------------------------------------------------------
% BAB V. HASIL PENELITIAN DAN PEMBAHASAN
% ---------------------------------------------------------------------
\chapter{HASIL PENELITIAN DAN PEMBAHASAN}

% Susun hasil berdasarkan tujuan atau rumusan masalah.
\section{Hasil Penelitian}
Sajikan hasil utama penelitian secara sistematis. Gunakan persamaan, tabel,
gambar, grafik, atau diagram apabila membantu menjelaskan hasil. Setiap
hasil yang disajikan perlu dirujuk dan dijelaskan dalam uraian.

\section{Pembahasan}
Interpretasikan hasil berdasarkan teori dan metode yang telah dijelaskan.
Bandingkan dengan hasil penelitian terdahulu apabila relevan. Bahas makna
fisis, kecenderungan, ketidakpastian, keterbatasan, dan implikasi hasil
sesuai ruang lingkup penelitian.

% Tambahkan section/subsection berdasarkan kelompok hasil yang diperoleh.
\section{Pembahasan Tambahan (Jika Diperlukan)}
Gunakan bagian ini apabila terdapat hasil atau analisis lain yang belum
tercakup pada bagian sebelumnya. Hapus bagian ini jika tidak diperlukan.

% ---------------------------------------------------------------------
% BAB VI. KESIMPULAN DAN SARAN
% ---------------------------------------------------------------------
\chapter{KESIMPULAN DAN SARAN}

% Kesimpulan menjawab tujuan/rumusan masalah berdasarkan hasil.
\section{Kesimpulan}
Tuliskan kesimpulan utama secara ringkas, jelas, dan langsung menjawab tujuan
atau rumusan masalah. Kesimpulan harus didasarkan pada hasil penelitian yang
telah dibahas, bukan memuat data atau analisis baru.

% Saran berisi tindak lanjut yang dapat dilakukan.
\section{Saran}
Tuliskan saran untuk pengembangan penelitian, perbaikan metode, atau pekerjaan
lanjutan. Pastikan saran berangkat dari hasil dan keterbatasan penelitian.

% =====================================================================
% D. BAB TAMBAHAN: PANDUAN PENGGUNAAN FITUR LATEX
% =====================================================================
% Bab ini adalah bahan pembelajaran di dalam template, bukan bagian standar
% isi TUGAS AKHIR. Nonaktifkan sebelum naskah akhir diserahkan dengan mengubah
% \showlatexguidetrue menjadi \showlatexguidefalse di bagian preamble.

\ifshowlatexguide
\chapter{PANDUAN FITUR LaTeX}

Bab ini memperkenalkan perintah LaTeX yang umum digunakan ketika menulis
naskah ilmiah. Contoh-contoh berikut dapat disalin dan disesuaikan pada bab
penelitian. Bab panduan ini harus dinonaktifkan atau dihapus sebelum naskah
akhir diserahkan apabila tidak termasuk struktur yang ditetapkan Fakultas.

\section{Struktur Bab, Subbab, dan Paragraf}
Gunakan perintah \verb|\chapter| untuk membuat bab, \verb|\section| untuk
subbab tingkat pertama, dan \verb|\subsection| untuk bagian yang lebih rinci.
Nomor dan format judul akan diatur oleh document class.

Contoh penulisan:
\begin{verbatim}
\chapter{JUDUL BAB}
\section{Judul Subbab}
\subsection{Judul Bagian}
Tuliskan uraian paragraf di sini.
\end{verbatim}

\section{Penulisan Persamaan}
Persamaan singkat dapat ditulis di dalam paragraf, misalnya $E=mc^2$.
Untuk persamaan yang perlu diberi nomor dan dirujuk kembali, gunakan
lingkungan \texttt{equation}.

\begin{equation}
E = mc^2,
\label{eq:contoh-einstein}
\end{equation}

Persamaan multibaris dapat ditulis dengan lingkungan \texttt{align}.
Paket \texttt{amsmath} telah dipanggil pada preamble.

\begin{align}
F &= ma, \label{eq:hukum-newton}\\
E &= mc^2. \label{eq:energi-massa}
\end{align}

Rujuk persamaan menggunakan \verb|\ref| atau \verb|\eqref|, misalnya
Persamaan~\eqref{eq:contoh-einstein}. Gunakan \verb|\label| pada persamaan
agar nomor diperbarui secara otomatis.

Contoh kode:
\begin{verbatim}
\begin{equation}
    E = mc^2
    \label{eq:energi}
\end{equation}

Persamaan~\eqref{eq:energi} menunjukkan ...
\end{verbatim}

\section{Memasukkan Gambar}
Simpan gambar pada folder \texttt{figure/} di dalam direktori proyek.
Gunakan nama file sederhana dan panggil gambar melalui path relatif.
Class template telah memuat paket \texttt{graphicx}.

Contoh berikut menampilkan gambar apabila file contoh tersedia. Jika belum,
template menampilkan kotak pengganti sehingga dokumen tetap dapat dikompilasi.

\begin{figure}[htbp]
    \centering
    \IfFileExists{figure/contoh-gambar.png}{%
        \includegraphics[width=0.70\textwidth]{figure/contoh-gambar.png}%
    }{%
        \fbox{\parbox[c][35mm][c]{0.70\textwidth}{\centering
        Contoh gambar belum tersedia}}%
    }
    \caption{Contoh penyisipan gambar.}
    \label{fig:contoh-gambar}
\end{figure}

Rujuk gambar di dalam paragraf menggunakan perintah \verb|\ref|, seperti
Gambar~\ref{fig:contoh-gambar}. Jangan menuliskan nomor gambar secara manual.

Contoh kode untuk gambar lain:
\begin{verbatim}
\begin{figure}[htbp]
    \centering
    \includegraphics[width=0.8\textwidth]{figure/nama-gambar.png}
    \caption{Keterangan gambar.}
    \label{fig:nama-gambar}
\end{figure}
\end{verbatim}

\section{Membuat Tabel}
Gunakan lingkungan \texttt{table} untuk memberi keterangan dan nomor tabel,
serta \texttt{tabular} untuk menyusun baris dan kolom. Contoh berikut
menggunakan paket \texttt{booktabs}.

\begin{table}[htbp]
    \centering
    \caption{Contoh penyajian data dalam tabel.}
    \label{tab:contoh-data}
    \begin{tabular}{lcc}
        \toprule
        Besaran & Nilai 1 & Nilai 2 \\
        \midrule
        Besaran A & 1.25 & 2.50 \\
        Besaran B & 3.75 & 4.00 \\
        \bottomrule
    \end{tabular}
\end{table}

Rujuk tabel menggunakan Tabel~\ref{tab:contoh-data}. Letakkan \verb|\caption|
dan \verb|\label| di dalam lingkungan \texttt{table}; tempatkan label
setelah caption agar nomor yang dirujuk sesuai.

\section{Membuat Daftar dan Poin}
Gunakan \texttt{itemize} untuk daftar tanpa nomor dan \texttt{enumerate}
untuk daftar bernomor.

\begin{itemize}
    \item Poin pertama.
    \item Poin kedua.
\end{itemize}

\begin{enumerate}
    \item Langkah pertama.
    \item Langkah kedua.
\end{enumerate}

\section{Sitasi dan Daftar Pustaka}
Class template memuat paket \texttt{natbib} dengan gaya bibliografi
\texttt{apalike}. Database referensi disimpan dalam \texttt{reference.bib}.
Setiap entri memiliki citation key untuk dipanggil dalam sitasi.

Dua perintah sitasi yang umum adalah:
\begin{itemize}
    \item \verb|\citet{key}| untuk sitasi naratif, misalnya nama penulis menjadi
    bagian dari kalimat.
    \item \verb|\citep{key}| untuk sitasi di dalam tanda kurung.
\end{itemize}

Ganti \texttt{key} dengan citation key yang benar-benar tersedia di dalam
\texttt{reference.bib}. Contoh bentuk kode:
\begin{verbatim}
Menurut \citet{key-referensi}, ...
Hasil tersebut juga telah dibahas sebelumnya \citep{key-referensi}.
\end{verbatim}

Daftar pustaka dibuat dari database BibTeX melalui perintah berikut pada
bagian akhir dokumen:
\begin{verbatim}
\bibliography{reference}
\end{verbatim}

Jangan menuliskan ekstensi \texttt{.bib} pada argumen \verb|\bibliography|.
Pastikan proses kompilasi menjalankan BibTeX ketika terdapat perubahan sitasi.

\section{Referensi Silang}
Gunakan \verb|\label| untuk memberi penanda pada bab, subbab, gambar, tabel,
dan persamaan. Panggil penanda tersebut menggunakan \verb|\ref| atau
\verb|\eqref| untuk persamaan.

Contoh:
\begin{verbatim}
\section{Model Penelitian}
\label{sec:model}

Rujuk Bagian~\ref{sec:model}.
\end{verbatim}

Kompilasi ulang diperlukan agar nomor referensi silang diperbarui.

\section{Catatan Kaki}
Catatan kaki dapat ditambahkan menggunakan perintah \verb|\footnote|.
Contoh kalimat dengan catatan kaki.\footnote{Tuliskan keterangan tambahan
yang relevan dan tidak perlu dimasukkan ke dalam uraian utama.}

\section{Lambang dan Singkatan}
Class template memuat paket \texttt{nomencl}. Tambahkan entri menggunakan
perintah \verb|\nomenclature|, kemudian aktifkan perintah \verb|\lambang|
pada bagian awal dokumen apabila daftar lambang digunakan.

Contoh:
\begin{verbatim}
\nomenclature{$c$}{Kecepatan cahaya di ruang hampa}
\nomenclature{$m$}{Massa benda}
\end{verbatim}

Pembuatan daftar lambang memerlukan tahapan pemrosesan indeks nomenklatur
selain kompilasi LaTeX biasa. Ikuti konfigurasi dan perintah kompilasi yang
berlaku pada instalasi LaTeX yang digunakan.

\section{Lampiran}
Gunakan \verb|\appendix| sebelum membuat bab lampiran. Setelah perintah ini,
penomoran bab akan mengikuti format lampiran yang ditentukan oleh class.

Contoh:
\begin{verbatim}
\appendix
\chapter{DATA TAMBAHAN}
Tuliskan data atau informasi pendukung di sini.
\end{verbatim}

\section{Kompilasi dan Pembaruan Dokumen}
Setelah menambahkan gambar, persamaan, tabel, atau sitasi, lakukan kompilasi
ulang melalui LaTeX Workshop. Beberapa elemen seperti daftar isi, referensi
silang, dan daftar pustaka memerlukan beberapa tahap kompilasi agar hasilnya
terbarui. Periksa panel Output apabila proses kompilasi menghasilkan pesan
kesalahan.
\fi

% =====================================================================
% E. DAFTAR PUSTAKA
% =====================================================================
% Class TUGAS AKHIR dari Fakultas pada template asli mengatur gaya apalike.
% File database yang digunakan pada starter project bernama reference.bib.
\bibliography{reference}

% =====================================================================
% F. LAMPIRAN
% =====================================================================
% Hapus contoh ini jika tidak ada lampiran. Tambahkan \chapter untuk setiap
% lampiran baru; seluruh lampiran tetap ditulis dalam main.tex.
\appendix
\chapter{JUDUL LAMPIRAN}
Tuliskan materi pendukung yang diperlukan, misalnya penurunan matematis,
data tambahan, spesifikasi alat, atau kode yang relevan. Sesuaikan jenis
lampiran dengan kebutuhan penelitian dan ketentuan panduan Fakultas.

% =====================================================================
% AKHIR DOKUMEN
% =====================================================================
\end{document}